\documentclass[11pt,a4paper]{article}

%========================
% PACKAGES
%========================
\usepackage[utf8]{inputenc}
\usepackage[T1]{fontenc}
\usepackage[french]{babel}
\usepackage[a4paper,margin=2cm]{geometry}

\usepackage{amsmath,amssymb,amsthm}
\usepackage{lmodern}
\usepackage{xcolor}
\usepackage{fancyhdr}
\usepackage{lastpage}
\usepackage{tikz}
\usepackage[most]{tcolorbox}
\usepackage{enumitem}
\usepackage{array}
\usepackage{tabularx}
\usepackage{multicol}

%========================
% COULEURS
%========================
\definecolor{bleuprepa}{RGB}{20,55,110}
\definecolor{vertprepa}{RGB}{20,105,70}
\definecolor{rougeprepa}{RGB}{160,35,35}
\definecolor{orangeprepa}{RGB}{180,100,20}
\definecolor{violetprepa}{RGB}{90,40,130}
\definecolor{grisclair}{RGB}{245,247,250}

%========================
% MISE EN PAGE
%========================
\setlength{\headheight}{14pt}
\pagestyle{fancy}
\fancyhf{}
\fancyhead[L]{Mathématiques -- Troisième}
\fancyhead[C]{Données, probabilités, proportionnalité, grandeurs}
\fancyhead[R]{Page \thepage/\pageref{LastPage}}
\fancyfoot[C]{Cours complet -- Cycle 4}

\setlength{\parindent}{0pt}
\setlength{\parskip}{0.4em}

%========================
% BOÎTES
%========================
\tcbset{
  enhanced,
  breakable,
  boxrule=0.7pt,
  arc=1mm,
  colback=white,
  fonttitle=\bfseries,
  before skip=8pt,
  after skip=8pt
}

\newtcolorbox{definitionbox}[1]{
  colframe=bleuprepa,
  colback=bleuprepa!4!white,
  title={Définition -- #1}
}

\newtcolorbox{propositionbox}[1]{
  colframe=vertprepa,
  colback=vertprepa!4!white,
  title={Propriété -- #1}
}

\newtcolorbox{methodbox}[1]{
  colframe=orangeprepa,
  colback=orangeprepa!5!white,
  title={Méthode -- #1}
}

\newtcolorbox{warningbox}[1]{
  colframe=rougeprepa,
  colback=rougeprepa!4!white,
  title={Erreur classique -- #1}
}

\newtcolorbox{retenirbox}[1]{
  colframe=violetprepa,
  colback=violetprepa!5!white,
  title={À retenir -- #1}
}

\newtcolorbox{examplebox}[1]{
  colframe=black!60,
  colback=grisclair,
  title={Exemple corrigé -- #1}
}

\newtcolorbox{exercisebox}[1]{
  colframe=black,
  colback=white,
  title={Exercice -- #1}
}

\newtcolorbox{correctionbox}[1]{
  colframe=vertprepa,
  colback=vertprepa!3!white,
  title={Correction -- #1}
}

%========================
% COMMANDES
%========================
\newcommand{\R}{\mathbb{R}}
\newcommand{\Prob}{\mathbb{P}}
\newcommand{\freq}{\mathrm{f}}
\newcommand{\eff}{\mathrm{n}}
\newcommand{\intervalle}[2]{\ensuremath{[#1\,;#2[}}

%========================
% DOCUMENT
%========================
\begin{document}

%========================
% PAGE DE TITRE
%========================
\begin{center}
    {\Huge \textbf{Cours complet de mathématiques}}\\[0.3cm]
    {\Large Classe de troisième}\\[0.3cm]
    {\Large \textbf{Données, probabilités, proportionnalité et grandeurs}}\\[0.5cm]
    \rule{0.9\linewidth}{0.5pt}\\[0.3cm]
    \textit{Lire, représenter, calculer, modéliser et interpréter}\\[0.3cm]
    \rule{0.9\linewidth}{0.5pt}
\end{center}

\vspace{0.5cm}

\begin{retenirbox}{Objectifs du chapitre}
À la fin de ce cours, l'élève doit savoir :
\begin{itemize}
    \item lire et représenter des données à l'aide de tableaux, diagrammes et histogrammes ;
    \item calculer des effectifs, fréquences, pourcentages et étendues ;
    \item interpréter une série statistique ;
    \item calculer des probabilités dans des expériences simples ou à deux épreuves ;
    \item faire le lien entre fréquence observée et probabilité théorique ;
    \item reconnaître et utiliser une situation de proportionnalité ;
    \item modéliser une proportionnalité par une fonction linéaire ;
    \item utiliser les coefficients multiplicateurs pour les évolutions en pourcentage ;
    \item calculer avec des grandeurs mesurables : longueurs, aires, volumes, vitesses, débits ;
    \item convertir correctement des unités simples ou composées ;
    \item vérifier la cohérence d'un résultat grâce aux unités.
\end{itemize}
\end{retenirbox}

\tableofcontents
\newpage

%================================================
\section{Interpréter, représenter et traiter des données}

\subsection{Vocabulaire statistique de base}

\begin{definitionbox}{Série statistique}
Une \textbf{série statistique} est un ensemble de données recueillies lors d'une enquête, d'une mesure ou d'une observation.

Chaque donnée correspond à une \textbf{valeur} ou à une \textbf{modalité}. Le nombre total de données s'appelle l'\textbf{effectif total}.
\end{definitionbox}

\begin{examplebox}{Pointures d'une classe}
On interroge $12$ élèves sur leur pointure :
\[
38,\ 39,\ 40,\ 42,\ 41,\ 40,\ 39,\ 38,\ 43,\ 41,\ 40,\ 42.
\]
Cette liste est une série statistique.

L'effectif total est :
\[
N=12.
\]
La valeur $40$ apparaît $3$ fois : son effectif est donc $3$.
\end{examplebox}

\begin{definitionbox}{Effectif}
L'\textbf{effectif} d'une valeur est le nombre de fois où cette valeur apparaît dans la série.
\end{definitionbox}

\begin{definitionbox}{Fréquence}
La \textbf{fréquence} d'une valeur est le quotient :
\[
\text{fréquence}=\frac{\text{effectif de la valeur}}{\text{effectif total}}.
\]
Si l'effectif d'une valeur est $n$ et l'effectif total est $N$, alors :
\[
f=\frac{n}{N}.
\]
La fréquence peut être écrite sous forme décimale, fractionnaire ou en pourcentage.
\end{definitionbox}

\begin{examplebox}{Calculer une fréquence}
Dans une classe de $25$ élèves, $10$ élèves viennent au collège en bus.

La fréquence des élèves venant en bus est :
\[
f=\frac{10}{25}=\frac{2}{5}=0{,}4.
\]
En pourcentage :
\[
0{,}4=40\%.
\]

Donc $40\%$ des élèves viennent en bus.
\end{examplebox}

\begin{methodbox}{Calculer une fréquence}
Pour calculer une fréquence :
\begin{enumerate}
    \item identifier l'effectif de la valeur étudiée ;
    \item identifier l'effectif total ;
    \item calculer le quotient :
    \[
    f=\frac{\text{effectif}}{\text{effectif total}} ;
    \]
    \item convertir en pourcentage si nécessaire.
\end{enumerate}
\end{methodbox}

\begin{warningbox}{Confondre effectif et fréquence}
L'effectif est un \textbf{nombre d'individus}.  
La fréquence est une \textbf{proportion}.  

Par exemple, si $8$ élèves sur $20$ ont réussi, $8$ est un effectif, tandis que :
\[
\frac{8}{20}=0{,}4=40\%
\]
est une fréquence.
\end{warningbox}

\subsection{Tableaux statistiques}

Un tableau statistique permet de ranger les données pour les rendre plus lisibles.

\begin{examplebox}{Tableau d'effectifs et de fréquences}
On relève les notes obtenues par $20$ élèves à un test noté sur $5$.

\[
\begin{array}{|c|c|c|}
\hline
\text{Note} & \text{Effectif} & \text{Fréquence} \\
\hline
1 & 2 & \frac{2}{20}=0{,}10=10\% \\
2 & 4 & \frac{4}{20}=0{,}20=20\% \\
3 & 6 & \frac{6}{20}=0{,}30=30\% \\
4 & 5 & \frac{5}{20}=0{,}25=25\% \\
5 & 3 & \frac{3}{20}=0{,}15=15\% \\
\hline
\text{Total} & 20 & 1=100\% \\
\hline
\end{array}
\]

On vérifie :
\[
2+4+6+5+3=20
\]
et
\[
10\%+20\%+30\%+25\%+15\%=100\%.
\]
\end{examplebox}

\begin{retenirbox}{Somme des fréquences}
Dans une série statistique complète, la somme des fréquences vaut toujours :
\[
1.
\]
En pourcentage, la somme vaut :
\[
100\%.
\]
\end{retenirbox}

\subsection{Diagramme en bâtons}

\begin{definitionbox}{Diagramme en bâtons}
Un \textbf{diagramme en bâtons} est une représentation graphique adaptée aux données discrètes.

Chaque valeur est représentée par un bâton dont la hauteur est proportionnelle à l'effectif ou à la fréquence.
\end{definitionbox}

\begin{examplebox}{Diagramme en bâtons}
On considère le tableau suivant :

\[
\begin{array}{|c|c|}
\hline
\text{Nombre de frères et sœurs} & \text{Effectif} \\
\hline
0 & 4 \\
1 & 8 \\
2 & 6 \\
3 & 2 \\
\hline
\end{array}
\]

Le diagramme en bâtons correspondant est :

\begin{center}
\begin{tikzpicture}[scale=0.8]
\draw[->] (-0.5,0) -- (4,0) node[right]{Nombre};
\draw[->] (0,-0.2) -- (0,9) node[above]{Effectif};

\foreach \y in {0,1,2,3,4,5,6,7,8}
{
  \draw (-0.1,\y) -- (0.1,\y);
  \node[left] at (-0.1,\y) {\small \y};
}

\foreach \x/\h in {0/4,1/8,2/6,3/2}
{
  \draw[line width=3pt,bleuprepa] (\x+0.5,0) -- (\x+0.5,\h);
  \node[below] at (\x+0.5,0) {\x};
}
\end{tikzpicture}
\end{center}
\end{examplebox}

\subsection{Diagramme circulaire}

\begin{definitionbox}{Diagramme circulaire}
Un \textbf{diagramme circulaire} représente les fréquences par des secteurs d'un disque.

L'angle d'un secteur est proportionnel à la fréquence de la valeur représentée.
\end{definitionbox}

\begin{propositionbox}{Calcul de l'angle d'un secteur}
Si une valeur a pour fréquence $f$, alors l'angle du secteur correspondant est :
\[
\theta=f\times 360^\circ.
\]

Si la fréquence est donnée en pourcentage $p\%$, alors :
\[
\theta=\frac{p}{100}\times 360^\circ.
\]
\end{propositionbox}

\begin{examplebox}{Construire un diagramme circulaire}
Dans un collège, on interroge $100$ élèves sur leur sport préféré.

\[
\begin{array}{|c|c|c|}
\hline
\text{Sport} & \text{Fréquence} & \text{Angle} \\
\hline
Football & 40\% & 0{,}40\times 360^\circ=144^\circ \\
Basket & 25\% & 0{,}25\times 360^\circ=90^\circ \\
Tennis & 15\% & 0{,}15\times 360^\circ=54^\circ \\
Autres & 20\% & 0{,}20\times 360^\circ=72^\circ \\
\hline
\end{array}
\]

On vérifie :
\[
144^\circ+90^\circ+54^\circ+72^\circ=360^\circ.
\]
\end{examplebox}

\subsection{Histogrammes}

\begin{definitionbox}{Classes}
Lorsque les données sont nombreuses ou continues, on peut les regrouper en \textbf{classes}.

Par exemple :
\[
\intervalle{18}{22},\quad \intervalle{22}{26},\quad \intervalle{26}{30}
\]
sont des classes d'âges.

La notation $\intervalle{18}{22}$ signifie :
\[
18 \leq x < 22.
\]
Le nombre $18$ est inclus, mais le nombre $22$ ne l'est pas.
\end{definitionbox}

\begin{definitionbox}{Histogramme pour classes de même amplitude}
Un \textbf{histogramme} représente des données regroupées en classes.

Lorsque les classes ont la même amplitude, chaque classe est représentée par un rectangle :
\begin{itemize}
    \item la largeur du rectangle correspond à l'amplitude de la classe ;
    \item la hauteur du rectangle est proportionnelle à l'effectif de la classe.
\end{itemize}
\end{definitionbox}

\begin{warningbox}{Histogramme et diagramme en bâtons}
Un diagramme en bâtons convient pour des valeurs séparées : $0$, $1$, $2$, $3$.

Un histogramme convient pour des intervalles :
\[
\intervalle{18}{22},\ \intervalle{22}{26},\ \intervalle{26}{30}.
\]

Dans un histogramme, les rectangles sont accolés.
\end{warningbox}

\begin{examplebox}{Représenter une enquête par un histogramme}
Une enquête a été réalisée auprès de $2500$ personnes à partir de la question :

\og À quel âge avez-vous trouvé un emploi correspondant à votre qualification ? \fg

\[
\begin{array}{|c|c|}
\hline
\text{Âge} & \text{Effectif} \\
\hline
\intervalle{18}{22} & 100 \\
\intervalle{22}{26} & 200 \\
\intervalle{26}{30} & 400 \\
\intervalle{30}{34} & 1100 \\
\intervalle{34}{38} & 700 \\
\hline
\end{array}
\]

Toutes les classes ont la même amplitude :
\[
22-18=26-22=30-26=34-30=38-34=4.
\]

On peut donc représenter les effectifs directement par les hauteurs des rectangles.

\begin{center}
\begin{tikzpicture}[xscale=0.35,yscale=0.006]
\draw[->] (17,0) -- (39,0) node[right]{Âge};
\draw[->] (18,0) -- (18,1250) node[above]{Effectif};

\foreach \x in {18,22,26,30,34,38}
{
  \draw (\x,0) -- (\x,-30);
  \node[below] at (\x,-30) {\small \x};
}

\foreach \y in {0,200,400,600,800,1000,1200}
{
  \draw (17.8,\y) -- (18.2,\y);
  \node[left] at (17.8,\y) {\small \y};
}

\draw[fill=blue!20,draw=bleuprepa] (18,0) rectangle (22,100);
\draw[fill=blue!20,draw=bleuprepa] (22,0) rectangle (26,200);
\draw[fill=blue!20,draw=bleuprepa] (26,0) rectangle (30,400);
\draw[fill=blue!20,draw=bleuprepa] (30,0) rectangle (34,1100);
\draw[fill=blue!20,draw=bleuprepa] (34,0) rectangle (38,700);

\node at (20,160) {\small 100};
\node at (24,260) {\small 200};
\node at (28,460) {\small 400};
\node at (32,1160) {\small 1100};
\node at (36,760) {\small 700};
\end{tikzpicture}
\end{center}
\end{examplebox}

\subsection{Étendue d'une série}

\begin{definitionbox}{Étendue}
L'\textbf{étendue} d'une série statistique est la différence entre la plus grande valeur et la plus petite valeur :
\[
\text{étendue}=\text{valeur maximale}-\text{valeur minimale}.
\]
\end{definitionbox}

\begin{examplebox}{Calculer une étendue avec des données brutes}
On considère la série :
\[
12,\ 15,\ 9,\ 18,\ 11,\ 20,\ 14.
\]

La plus petite valeur est $9$ et la plus grande est $20$.

Donc :
\[
\text{étendue}=20-9=11.
\]

L'étendue est $11$.
\end{examplebox}

\begin{examplebox}{Étendue avec des classes}
On considère des âges regroupés ainsi :
\[
\intervalle{18}{22},\quad \intervalle{22}{26},\quad \intervalle{26}{30},\quad \intervalle{30}{34}.
\]

Dans ce cas, on ne connaît pas les valeurs exactes. On peut seulement dire que les âges vont de $18$ jusqu'à moins de $34$.

Une étendue approximative ou maximale raisonnable est :
\[
34-18=16.
\]

Mais il faut dire clairement que l'on travaille avec des données regroupées.
\end{examplebox}

\begin{warningbox}{Étendue et données regroupées}
Lorsque les données sont regroupées en classes, on ne connaît pas les valeurs individuelles exactes. L'étendue calculée à partir des bornes des classes est donc une estimation ou une interprétation selon le contexte.
\end{warningbox}

%================================================
\subsection{Exercices progressifs -- Données}

\begin{exercisebox}{Effectifs et fréquences}
Dans une classe de $30$ élèves, on relève le moyen de transport utilisé pour venir au collège.

\[
\begin{array}{|c|c|}
\hline
\text{Transport} & \text{Effectif} \\
\hline
À pied & 6 \\
Bus & 12 \\
Voiture & 9 \\
Vélo & 3 \\
\hline
\end{array}
\]

\begin{enumerate}
    \item Calculer la fréquence des élèves venant en bus.
    \item Calculer la fréquence des élèves venant à pied.
    \item Donner ces fréquences en pourcentage.
    \item Vérifier que la somme des fréquences vaut $1$.
\end{enumerate}
\end{exercisebox}

\begin{exercisebox}{Étendue}
Voici les températures maximales relevées pendant une semaine :
\[
18,\ 21,\ 20,\ 16,\ 23,\ 19,\ 25.
\]
\begin{enumerate}
    \item Donner la température minimale.
    \item Donner la température maximale.
    \item Calculer l'étendue.
    \item Interpréter le résultat par une phrase.
\end{enumerate}
\end{exercisebox}

\begin{exercisebox}{Histogramme}
On a interrogé $80$ personnes sur le temps quotidien passé à lire.

\[
\begin{array}{|c|c|}
\hline
\text{Temps en minutes} & \text{Effectif} \\
\hline
\intervalle{0}{15} & 12 \\
\intervalle{15}{30} & 20 \\
\intervalle{30}{45} & 28 \\
\intervalle{45}{60} & 16 \\
\intervalle{60}{75} & 4 \\
\hline
\end{array}
\]

\begin{enumerate}
    \item Vérifier que l'effectif total est bien $80$.
    \item Quelle est l'amplitude de chaque classe ?
    \item Représenter les données par un histogramme.
    \item Calculer la fréquence des personnes lisant entre $30$ et $45$ minutes.
\end{enumerate}
\end{exercisebox}

\newpage

%================================================
\section{Comprendre et utiliser des probabilités élémentaires}

\subsection{Expérience aléatoire}

\begin{definitionbox}{Expérience aléatoire}
Une \textbf{expérience aléatoire} est une expérience dont on connaît les résultats possibles, mais dont on ne peut pas prévoir avec certitude le résultat qui va se produire.
\end{definitionbox}

\begin{examplebox}{Expériences aléatoires}
\begin{itemize}
    \item Lancer une pièce est une expérience aléatoire : les résultats possibles sont pile et face.
    \item Lancer un dé équilibré est une expérience aléatoire : les résultats possibles sont $1$, $2$, $3$, $4$, $5$, $6$.
    \item Tirer une carte au hasard dans un jeu est une expérience aléatoire.
\end{itemize}
\end{examplebox}

\begin{definitionbox}{Issue}
Une \textbf{issue} est un résultat possible d'une expérience aléatoire.
\end{definitionbox}

\begin{definitionbox}{Univers}
L'\textbf{univers} d'une expérience aléatoire est l'ensemble de toutes les issues possibles. On le note souvent $\Omega$.
\end{definitionbox}

\begin{examplebox}{Univers d'un lancer de dé}
Lorsqu'on lance un dé à six faces, l'univers est :
\[
\Omega=\{1,2,3,4,5,6\}.
\]
\end{examplebox}

\subsection{Événements}

\begin{definitionbox}{Événement}
Un \textbf{événement} est un ensemble d'issues.

Dire qu'un événement est réalisé signifie que le résultat obtenu appartient à cet événement.
\end{definitionbox}

\begin{examplebox}{Événement avec un dé}
On lance un dé équilibré.

L'événement $A$ : \og obtenir un nombre pair \fg{} est :
\[
A=\{2,4,6\}.
\]

L'événement $B$ : \og obtenir un nombre supérieur ou égal à $5$ \fg{} est :
\[
B=\{5,6\}.
\]
\end{examplebox}

\begin{definitionbox}{Événement contraire}
L'\textbf{événement contraire} de $A$, noté $\overline{A}$, est l'événement qui se réalise lorsque $A$ ne se réalise pas.
\end{definitionbox}

\begin{examplebox}{Événement contraire}
On lance un dé.

Si $A$ est l'événement \og obtenir un nombre pair \fg, alors :
\[
A=\{2,4,6\}.
\]
Son contraire est :
\[
\overline{A}=\{1,3,5\}.
\]
\end{examplebox}

\subsection{Probabilité}

\begin{definitionbox}{Probabilité}
La \textbf{probabilité} d'un événement mesure la chance qu'a cet événement de se produire.

Une probabilité est un nombre compris entre $0$ et $1$.
\[
0\leq \Prob(A)\leq 1.
\]
\end{definitionbox}

\begin{retenirbox}{Cas particuliers}
\begin{itemize}
    \item Si $\Prob(A)=0$, l'événement est impossible.
    \item Si $\Prob(A)=1$, l'événement est certain.
    \item Plus $\Prob(A)$ est proche de $1$, plus l'événement a de chances de se produire.
    \item Plus $\Prob(A)$ est proche de $0$, moins l'événement a de chances de se produire.
\end{itemize}
\end{retenirbox}

\begin{propositionbox}{Situation d'équiprobabilité}
Lorsque toutes les issues ont la même probabilité, on dit qu'il y a \textbf{équiprobabilité}.

Dans ce cas :
\[
\Prob(A)=\frac{\text{nombre d'issues favorables à }A}{\text{nombre total d'issues}}.
\]
\end{propositionbox}

\begin{examplebox}{Probabilité avec un dé équilibré}
On lance un dé équilibré. On cherche la probabilité d'obtenir un nombre pair.

Les issues possibles sont :
\[
\Omega=\{1,2,3,4,5,6\}.
\]

Les issues favorables sont :
\[
\{2,4,6\}.
\]

Il y a donc $3$ issues favorables sur $6$ issues possibles.

Ainsi :
\[
\Prob(A)=\frac{3}{6}=\frac{1}{2}.
\]
\end{examplebox}

\begin{propositionbox}{Probabilité de l'événement contraire}
Pour tout événement $A$ :
\[
\Prob(\overline{A})=1-\Prob(A).
\]

De façon équivalente :
\[
\Prob(A)=1-\Prob(\overline{A}).
\]
\end{propositionbox}

\begin{examplebox}{Utiliser l'événement contraire}
On lance un dé. On cherche la probabilité d'obtenir au moins $2$.

L'événement contraire est : obtenir strictement moins que $2$, c'est-à-dire obtenir $1$.

Donc :
\[
\Prob(\text{obtenir }1)=\frac{1}{6}.
\]

Ainsi :
\[
\Prob(\text{obtenir au moins }2)=1-\frac{1}{6}=\frac{5}{6}.
\]
\end{examplebox}

\subsection{Expériences à deux épreuves}

\begin{examplebox}{Deux enfants}
On suppose que, pour un couple, la probabilité d'avoir une fille ou un garçon est la même. Un couple souhaite avoir deux enfants.

On note $F$ pour fille et $G$ pour garçon.

Les issues possibles sont :
\[
(F,F),\quad (F,G),\quad (G,F),\quad (G,G).
\]

Il y a $4$ issues équiprobables.

La probabilité d'avoir deux garçons est :
\[
\Prob(G,G)=\frac{1}{4}.
\]

La probabilité d'avoir au moins une fille peut être calculée avec l'événement contraire.

L'événement contraire de \og avoir au moins une fille \fg{} est \og avoir deux garçons \fg.

Donc :
\[
\Prob(\text{au moins une fille})=1-\Prob(G,G)=1-\frac{1}{4}=\frac{3}{4}.
\]
\end{examplebox}

\begin{methodbox}{Étudier une expérience à deux épreuves}
Pour une expérience à deux épreuves :
\begin{enumerate}
    \item lister les résultats possibles de la première épreuve ;
    \item lister les résultats possibles de la seconde épreuve ;
    \item construire un tableau à double entrée ou un arbre ;
    \item compter toutes les issues possibles ;
    \item compter les issues favorables ;
    \item calculer la probabilité.
\end{enumerate}
\end{methodbox}

\begin{examplebox}{Tirage avec remise}
On tire deux fois de suite, avec remise, une boule dans une urne contenant :
\begin{itemize}
    \item une boule bleue notée $B$ ;
    \item deux boules violettes notées $V_1$ et $V_2$.
\end{itemize}

Comme il y a remise, la composition de l'urne est la même au deuxième tirage.

On peut construire le tableau :

\[
\begin{array}{|c|c|c|c|}
\hline
 & B & V_1 & V_2 \\
\hline
B & (B,B) & (B,V_1) & (B,V_2) \\
\hline
V_1 & (V_1,B) & (V_1,V_1) & (V_1,V_2) \\
\hline
V_2 & (V_2,B) & (V_2,V_1) & (V_2,V_2) \\
\hline
\end{array}
\]

Il y a $9$ issues équiprobables.

Tirer deux boules violettes correspond aux issues :
\[
(V_1,V_1),\quad (V_1,V_2),\quad (V_2,V_1),\quad (V_2,V_2).
\]

Il y a donc $4$ issues favorables.

Ainsi :
\[
\Prob(\text{deux violettes})=\frac{4}{9}.
\]
\end{examplebox}

\subsection{Fréquences et probabilités}

\begin{definitionbox}{Fréquence observée}
Lorsqu'on répète une expérience aléatoire un grand nombre de fois, la \textbf{fréquence observée} d'un événement est :
\[
f=\frac{\text{nombre de réalisations de l'événement}}{\text{nombre total d'expériences}}.
\]
\end{definitionbox}

\begin{retenirbox}{Stabilisation des fréquences}
Lorsqu'on répète une expérience aléatoire un très grand nombre de fois, la fréquence observée d'un événement tend à se stabiliser autour de sa probabilité théorique.

Par exemple, pour un dé équilibré, la probabilité d'obtenir $6$ est :
\[
\frac{1}{6}\approx 0{,}167.
\]

Sur un petit nombre de lancers, la fréquence peut être éloignée de $0{,}167$.

Sur un grand nombre de lancers, elle se rapproche généralement de $0{,}167$.
\end{retenirbox}

\begin{warningbox}{Probabilité et certitude}
Dire qu'un événement a une probabilité $\frac{1}{6}$ ne signifie pas qu'il se produit exactement une fois tous les six essais.

Cela signifie qu'à long terme, sa fréquence tend à se rapprocher de $\frac{1}{6}$.
\end{warningbox}

%================================================
\subsection{Exercices progressifs -- Probabilités}

\begin{exercisebox}{Dé équilibré}
On lance un dé équilibré à six faces.

\begin{enumerate}
    \item Donner l'univers de l'expérience.
    \item Calculer la probabilité d'obtenir $4$.
    \item Calculer la probabilité d'obtenir un nombre pair.
    \item Calculer la probabilité d'obtenir un nombre strictement supérieur à $4$.
    \item Calculer la probabilité de ne pas obtenir $6$.
\end{enumerate}
\end{exercisebox}

\begin{exercisebox}{Deux pièces}
On lance deux pièces équilibrées.

\begin{enumerate}
    \item Lister toutes les issues possibles.
    \item Calculer la probabilité d'obtenir deux piles.
    \item Calculer la probabilité d'obtenir exactement une fois pile.
    \item Calculer la probabilité d'obtenir au moins une fois face.
\end{enumerate}
\end{exercisebox}

\begin{exercisebox}{Urne}
Une urne contient $3$ boules rouges, $2$ boules bleues et $5$ boules vertes. On tire une boule au hasard.

\begin{enumerate}
    \item Combien y a-t-il de boules au total ?
    \item Calculer la probabilité de tirer une boule rouge.
    \item Calculer la probabilité de tirer une boule bleue ou verte.
    \item Calculer la probabilité de ne pas tirer une boule verte.
\end{enumerate}
\end{exercisebox}

\newpage

%================================================
\section{Résoudre des problèmes de proportionnalité}

\subsection{Reconnaître une situation de proportionnalité}

\begin{definitionbox}{Proportionnalité}
Deux grandeurs sont \textbf{proportionnelles} lorsque l'on obtient toujours les valeurs de l'une en multipliant les valeurs de l'autre par un même nombre non nul.

Ce nombre s'appelle le \textbf{coefficient de proportionnalité}.
\end{definitionbox}

\begin{examplebox}{Situation proportionnelle}
Un mobile avance à vitesse constante de $5$ mètres par seconde.

\[
\begin{array}{|c|c|c|c|c|}
\hline
\text{Temps en s} & 1 & 2 & 3 & 10 \\
\hline
\text{Distance en m} & 5 & 10 & 15 & 50 \\
\hline
\end{array}
\]

La distance est toujours égale au temps multiplié par $5$.

Il s'agit donc d'une situation de proportionnalité.
\end{examplebox}

\begin{warningbox}{Toutes les situations ne sont pas proportionnelles}
Une situation n'est pas proportionnelle simplement parce que les deux grandeurs augmentent en même temps.

Par exemple, l'âge d'un enfant et sa taille ne sont pas proportionnels.
\end{warningbox}

\subsection{Fonction linéaire}

\begin{definitionbox}{Fonction linéaire}
Une fonction linéaire est une fonction de la forme :
\[
f(x)=ax
\]
où $a$ est un nombre fixé.

Le nombre $a$ est appelé \textbf{coefficient}.
\end{definitionbox}

\begin{propositionbox}{Proportionnalité et fonction linéaire}
Une situation de proportionnalité peut être modélisée par une fonction linéaire :
\[
f(x)=ax.
\]

Le coefficient $a$ est le coefficient de proportionnalité.
\end{propositionbox}

\begin{examplebox}{Mobile à vitesse constante}
Un mobile se déplace à $5$ m/s.

Si $x$ désigne le temps en secondes et $d(x)$ la distance parcourue en mètres, alors :
\[
d(x)=5x.
\]

Ainsi :
\[
d(8)=5\times 8=40.
\]

En $8$ secondes, le mobile parcourt $40$ mètres.
\end{examplebox}

\subsection{Pourcentages d'évolution}

\begin{definitionbox}{Augmentation en pourcentage}
Augmenter une quantité de $p\%$, c'est la multiplier par :
\[
1+\frac{p}{100}.
\]
\end{definitionbox}

\begin{definitionbox}{Diminution en pourcentage}
Diminuer une quantité de $p\%$, c'est la multiplier par :
\[
1-\frac{p}{100}.
\]
\end{definitionbox}

\begin{examplebox}{Augmentation de $5\%$}
Un prix de $80$ euros augmente de $5\%$.

Le coefficient multiplicateur est :
\[
1+\frac{5}{100}=1{,}05.
\]

Le nouveau prix est :
\[
80\times 1{,}05=84.
\]

Le nouveau prix est $84$ euros.
\end{examplebox}

\begin{examplebox}{Diminution de $20\%$}
Un article coûte $50$ euros. Il diminue de $20\%$.

Le coefficient multiplicateur est :
\[
1-\frac{20}{100}=0{,}8.
\]

Le nouveau prix est :
\[
50\times 0{,}8=40.
\]

Le nouveau prix est $40$ euros.
\end{examplebox}

\begin{warningbox}{Ajouter ou enlever le pourcentage}
Augmenter de $5\%$ ne signifie pas ajouter $5$.

Cela signifie ajouter $5\%$ de la valeur initiale.

Par exemple, $5\%$ de $80$ vaut :
\[
0{,}05\times 80=4.
\]
Donc $80$ augmenté de $5\%$ vaut :
\[
80+4=84.
\]
\end{warningbox}

\begin{methodbox}{Résoudre un problème d'évolution}
Pour résoudre un problème avec pourcentage d'évolution :
\begin{enumerate}
    \item identifier s'il s'agit d'une augmentation ou d'une diminution ;
    \item écrire le coefficient multiplicateur ;
    \item multiplier la valeur initiale par ce coefficient ;
    \item vérifier que le résultat est cohérent.
\end{enumerate}
\end{methodbox}

\subsection{Proportionnalité en géométrie}

\begin{retenirbox}{Géométrie}
La proportionnalité intervient dans plusieurs situations géométriques :
\begin{itemize}
    \item théorème de Thalès ;
    \item triangles semblables ;
    \item agrandissements et réductions ;
    \item homothéties.
\end{itemize}
\end{retenirbox}

\begin{examplebox}{Triangles semblables}
Deux triangles sont semblables. Les longueurs du deuxième triangle sont obtenues en multipliant celles du premier par $1{,}5$.

Si un côté du premier triangle mesure $8$ cm, alors le côté correspondant du second mesure :
\[
8\times 1{,}5=12\text{ cm}.
\]
\end{examplebox}

\begin{examplebox}{Thalès}
Dans une configuration de Thalès, on obtient :
\[
\frac{AM}{AB}=\frac{AN}{AC}.
\]

Si :
\[
AM=3,\quad AB=9,\quad AC=12,
\]
alors :
\[
\frac{3}{9}=\frac{AN}{12}.
\]

Comme :
\[
\frac{3}{9}=\frac{1}{3},
\]
on a :
\[
\frac{AN}{12}=\frac{1}{3}.
\]

Donc :
\[
AN=12\times \frac{1}{3}=4.
\]
\end{examplebox}

%================================================
\subsection{Exercices progressifs -- Proportionnalité}

\begin{exercisebox}{Vitesse constante}
Un cycliste roule à vitesse constante. Il parcourt $18$ km en $1$ heure.

\begin{enumerate}
    \item Quelle distance parcourt-il en $2$ heures ?
    \item Quelle distance parcourt-il en $3{,}5$ heures ?
    \item Modéliser la distance $d(t)$ parcourue en fonction du temps $t$ exprimé en heures.
    \item Calculer $d(5)$.
\end{enumerate}
\end{exercisebox}

\begin{exercisebox}{Prix et pourcentage}
Un ordinateur coûte $650$ euros. Son prix baisse de $12\%$.

\begin{enumerate}
    \item Donner le coefficient multiplicateur correspondant.
    \item Calculer le nouveau prix.
    \item Calculer le montant de la réduction.
\end{enumerate}
\end{exercisebox}

\begin{exercisebox}{Agrandissement}
Une figure est agrandie avec un coefficient $2{,}4$.

\begin{enumerate}
    \item Une longueur de $5$ cm devient quelle longueur ?
    \item Une longueur de $12$ cm devient quelle longueur ?
    \item Si une longueur de la figure agrandie mesure $36$ cm, quelle était la longueur initiale ?
\end{enumerate}
\end{exercisebox}

\newpage

%================================================
\section{Calculer avec des grandeurs mesurables}

\subsection{Grandeurs simples et grandeurs composées}

\begin{definitionbox}{Grandeur}
Une \textbf{grandeur} est une quantité que l'on peut mesurer.

Exemples :
\begin{itemize}
    \item une longueur ;
    \item une aire ;
    \item un volume ;
    \item une durée ;
    \item une masse ;
    \item une vitesse ;
    \item un débit.
\end{itemize}
\end{definitionbox}

\begin{definitionbox}{Grandeur composée}
Une \textbf{grandeur composée} est une grandeur qui dépend de plusieurs grandeurs simples.

Exemples :
\[
\text{vitesse}=\frac{\text{distance}}{\text{temps}},
\]
\[
\text{débit}=\frac{\text{volume}}{\text{temps}}.
\]
\end{definitionbox}

\subsection{Unités usuelles}

\begin{retenirbox}{Unités}
\begin{center}
\begin{tabularx}{0.9\linewidth}{|X|X|}
\hline
\textbf{Grandeur} & \textbf{Unités usuelles} \\
\hline
Longueur & mm, cm, m, km \\
\hline
Aire & mm$^2$, cm$^2$, m$^2$, km$^2$ \\
\hline
Volume & mm$^3$, cm$^3$, m$^3$, L \\
\hline
Durée & s, min, h \\
\hline
Vitesse & m/s, km/h \\
\hline
Débit & m$^3$/s, L/min, L/h \\
\hline
\end{tabularx}
\end{center}
\end{retenirbox}

\subsection{Conversions}

\begin{retenirbox}{Conversions de base}
\[
1\text{ km}=1000\text{ m}
\]
\[
1\text{ m}=100\text{ cm}
\]
\[
1\text{ h}=60\text{ min}=3600\text{ s}
\]
\[
1\text{ L}=1\text{ dm}^3
\]
\[
1\text{ m}^3=1000\text{ L}
\]
\end{retenirbox}

\begin{methodbox}{Convertir une vitesse de km/h en m/s}
Pour convertir une vitesse de km/h en m/s :
\[
1\text{ km}=1000\text{ m}
\]
et
\[
1\text{ h}=3600\text{ s}.
\]

Donc :
\[
1\text{ km/h}=\frac{1000}{3600}\text{ m/s}=\frac{1}{3{,}6}\text{ m/s}.
\]

Ainsi, pour passer de km/h à m/s, on divise par $3{,}6$.
\end{methodbox}

\begin{examplebox}{Distance de réaction}
Un conducteur met $1$ seconde avant de commencer à freiner lorsqu'il voit un obstacle. Quelle distance parcourt-il pendant cette durée s'il roule à $80$ km/h ?

On convertit :
\[
80\text{ km/h}=\frac{80}{3{,}6}\text{ m/s}\approx 22{,}2\text{ m/s}.
\]

En $1$ seconde, il parcourt donc environ :
\[
22{,}2\times 1=22{,}2\text{ m}.
\]

Il parcourt environ $22$ mètres avant de commencer à freiner.
\end{examplebox}

\subsection{Volumes usuels}

\begin{retenirbox}{Formules de volume}
\[
\text{Volume d'un pavé droit}=L\times l\times h
\]
\[
\text{Volume d'un prisme droit}=\text{aire de la base}\times \text{hauteur}
\]
\[
\text{Volume d'un cylindre}=\pi r^2 h
\]
\[
\text{Volume d'une pyramide}=\frac{1}{3}\times \text{aire de la base}\times \text{hauteur}
\]
\[
\text{Volume d'un cône}=\frac{1}{3}\pi r^2 h
\]
\[
\text{Volume d'une boule}=\frac{4}{3}\pi r^3
\]
\end{retenirbox}

\begin{examplebox}{Volume d'une boule}
Calculer le volume d'une boule de rayon $3$ cm.

\[
V=\frac{4}{3}\pi r^3.
\]

Donc :
\[
V=\frac{4}{3}\pi \times 3^3.
\]

Or :
\[
3^3=27.
\]

Ainsi :
\[
V=\frac{4}{3}\pi \times 27=36\pi.
\]

Le volume exact est :
\[
36\pi\text{ cm}^3.
\]

Une valeur approchée est :
\[
36\pi\approx 113{,}1\text{ cm}^3.
\]
\end{examplebox}

\begin{examplebox}{Cylindre surmonté d'une demi-boule}
On considère un cylindre de rayon $4$ cm et de hauteur $10$ cm surmonté d'une demi-boule de même rayon.

Volume du cylindre :
\[
V_1=\pi r^2h=\pi\times 4^2\times 10=160\pi.
\]

Volume de la boule complète :
\[
V_{\text{boule}}=\frac{4}{3}\pi r^3=\frac{4}{3}\pi\times 4^3
=\frac{256}{3}\pi.
\]

Volume de la demi-boule :
\[
V_2=\frac{1}{2}\times \frac{256}{3}\pi=\frac{128}{3}\pi.
\]

Volume total :
\[
V=V_1+V_2=160\pi+\frac{128}{3}\pi.
\]

On met au même dénominateur :
\[
160\pi=\frac{480}{3}\pi.
\]

Donc :
\[
V=\frac{480}{3}\pi+\frac{128}{3}\pi=\frac{608}{3}\pi.
\]

Le volume total est :
\[
\frac{608}{3}\pi\text{ cm}^3.
\]
\end{examplebox}

\subsection{Débits}

\begin{definitionbox}{Débit}
Le \textbf{débit} est le volume écoulé par unité de temps.

\[
\text{débit}=\frac{\text{volume}}{\text{durée}}.
\]

Donc :
\[
\text{volume}=\text{débit}\times \text{durée}.
\]
\end{definitionbox}

\begin{examplebox}{Débit de la Seine}
Le débit moyen de la Seine sous le pont de l'Alma est $328$ m$^3$/s.

Combien de litres d'eau sont passés sous ce pont en $3$ minutes ?

On convertit d'abord la durée :
\[
3\text{ min}=3\times 60=180\text{ s}.
\]

Le volume écoulé est :
\[
V=328\times 180=59040\text{ m}^3.
\]

Or :
\[
1\text{ m}^3=1000\text{ L}.
\]

Donc :
\[
59040\text{ m}^3=59040\times 1000=59040000\text{ L}.
\]

Il est passé :
\[
59\,040\,000\text{ L}
\]
d'eau sous le pont en $3$ minutes.
\end{examplebox}

\begin{warningbox}{Oublier les unités}
Dans un calcul de grandeurs, les unités sont indispensables.

Écrire seulement :
\[
328\times 180=59040
\]
ne suffit pas.

Il faut écrire :
\[
328\text{ m}^3/\text{s}\times 180\text{ s}=59040\text{ m}^3.
\]
\end{warningbox}

%================================================
\subsection{Exercices progressifs -- Grandeurs}

\begin{exercisebox}{Volume d'une boule}
Calculer le volume exact puis une valeur approchée au dixième d'une boule de rayon $6$ cm.
\end{exercisebox}

\begin{exercisebox}{Vitesse et distance}
Une voiture roule à $90$ km/h.

\begin{enumerate}
    \item Convertir cette vitesse en m/s.
    \item Quelle distance parcourt-elle en $2$ secondes ?
    \item Quelle distance parcourt-elle en $5$ secondes ?
\end{enumerate}
\end{exercisebox}

\begin{exercisebox}{Débit}
Un robinet a un débit de $15$ L/min.

\begin{enumerate}
    \item Quel volume d'eau s'écoule en $8$ minutes ?
    \item Convertir ce volume en m$^3$.
    \item Combien de temps faut-il pour remplir une cuve de $450$ L ?
\end{enumerate}
\end{exercisebox}

\newpage

%================================================
\section{Corrections détaillées des exercices}

\subsection{Corrections -- Données}

\begin{correctionbox}{Effectifs et fréquences}
L'effectif total est :
\[
N=30.
\]

\begin{enumerate}
    \item Pour les élèves venant en bus :
    \[
    f=\frac{12}{30}=\frac{2}{5}=0{,}4.
    \]
    \item Pour les élèves venant à pied :
    \[
    f=\frac{6}{30}=\frac{1}{5}=0{,}2.
    \]
    \item En pourcentage :
    \[
    0{,}4=40\%,\qquad 0{,}2=20\%.
    \]
    \item Les fréquences sont :
    \[
    \frac{6}{30},\quad \frac{12}{30},\quad \frac{9}{30},\quad \frac{3}{30}.
    \]
    Leur somme vaut :
    \[
    \frac{6+12+9+3}{30}=\frac{30}{30}=1.
    \]
\end{enumerate}
\end{correctionbox}

\begin{correctionbox}{Étendue}
La série est :
\[
18,\ 21,\ 20,\ 16,\ 23,\ 19,\ 25.
\]

La température minimale est :
\[
16.
\]

La température maximale est :
\[
25.
\]

L'étendue est :
\[
25-16=9.
\]

Interprétation : durant cette semaine, l'écart entre la température maximale la plus basse et la température maximale la plus haute est de $9^\circ$C.
\end{correctionbox}

\begin{correctionbox}{Histogramme}
\begin{enumerate}
    \item On calcule l'effectif total :
    \[
    12+20+28+16+4=80.
    \]
    L'effectif total est bien $80$.
    
    \item L'amplitude de chaque classe est :
    \[
    15-0=30-15=45-30=60-45=75-60=15.
    \]
    Chaque classe a une amplitude de $15$ minutes.
    
    \item Les classes ont même amplitude. On peut donc construire un histogramme dont les hauteurs sont les effectifs.
    
    \item La fréquence des personnes lisant entre $30$ et $45$ minutes est :
    \[
    f=\frac{28}{80}=\frac{7}{20}=0{,}35=35\%.
    \]
\end{enumerate}
\end{correctionbox}

\subsection{Corrections -- Probabilités}

\begin{correctionbox}{Dé équilibré}
On lance un dé équilibré.

\begin{enumerate}
    \item L'univers est :
    \[
    \Omega=\{1,2,3,4,5,6\}.
    \]
    
    \item La probabilité d'obtenir $4$ est :
    \[
    \Prob(4)=\frac{1}{6}.
    \]
    
    \item Les nombres pairs sont :
    \[
    2,\ 4,\ 6.
    \]
    Il y a donc $3$ issues favorables sur $6$ :
    \[
    \Prob(\text{pair})=\frac{3}{6}=\frac{1}{2}.
    \]
    
    \item Les nombres strictement supérieurs à $4$ sont :
    \[
    5,\ 6.
    \]
    Donc :
    \[
    \Prob(>4)=\frac{2}{6}=\frac{1}{3}.
    \]
    
    \item L'événement contraire de \og obtenir $6$ \fg{} est \og ne pas obtenir $6$ \fg.
    \[
    \Prob(\text{ne pas obtenir }6)=1-\frac{1}{6}=\frac{5}{6}.
    \]
\end{enumerate}
\end{correctionbox}

\begin{correctionbox}{Deux pièces}
On note $P$ pour pile et $F$ pour face.

\begin{enumerate}
    \item Les issues possibles sont :
    \[
    (P,P),\quad (P,F),\quad (F,P),\quad (F,F).
    \]
    
    \item L'issue \og deux piles \fg{} est seulement :
    \[
    (P,P).
    \]
    Donc :
    \[
    \Prob(\text{deux piles})=\frac{1}{4}.
    \]
    
    \item Exactement une fois pile correspond à :
    \[
    (P,F),\quad (F,P).
    \]
    Donc :
    \[
    \Prob(\text{exactement une fois pile})=\frac{2}{4}=\frac{1}{2}.
    \]
    
    \item L'événement contraire de \og au moins une fois face \fg{} est \og aucune face \fg, c'est-à-dire \og deux piles \fg.
    Donc :
    \[
    \Prob(\text{au moins une face})=1-\frac{1}{4}=\frac{3}{4}.
    \]
\end{enumerate}
\end{correctionbox}

\begin{correctionbox}{Urne}
L'urne contient :
\[
3+2+5=10
\]
boules.

\begin{enumerate}
    \item Il y a $10$ boules au total.
    
    \item La probabilité de tirer une boule rouge est :
    \[
    \frac{3}{10}.
    \]
    
    \item Il y a :
    \[
    2+5=7
    \]
    boules bleues ou vertes.
    Donc :
    \[
    \Prob(\text{bleue ou verte})=\frac{7}{10}.
    \]
    
    \item Ne pas tirer une boule verte signifie tirer une boule rouge ou bleue.
    Il y a :
    \[
    3+2=5
    \]
    boules favorables.
    Donc :
    \[
    \Prob(\text{pas verte})=\frac{5}{10}=\frac{1}{2}.
    \]
\end{enumerate}
\end{correctionbox}

\subsection{Corrections -- Proportionnalité}

\begin{correctionbox}{Vitesse constante}
Le cycliste parcourt $18$ km en $1$ heure.

\begin{enumerate}
    \item En $2$ heures :
    \[
    18\times 2=36.
    \]
    Il parcourt $36$ km.
    
    \item En $3{,}5$ heures :
    \[
    18\times 3{,}5=63.
    \]
    Il parcourt $63$ km.
    
    \item Si $t$ désigne le temps en heures, alors :
    \[
    d(t)=18t.
    \]
    
    \item
    \[
    d(5)=18\times 5=90.
    \]
    En $5$ heures, il parcourt $90$ km.
\end{enumerate}
\end{correctionbox}

\begin{correctionbox}{Prix et pourcentage}
Le prix initial est $650$ euros. Il baisse de $12\%$.

\begin{enumerate}
    \item Le coefficient multiplicateur est :
    \[
    1-\frac{12}{100}=1-0{,}12=0{,}88.
    \]
    
    \item Le nouveau prix est :
    \[
    650\times 0{,}88=572.
    \]
    Le nouveau prix est $572$ euros.
    
    \item Le montant de la réduction est :
    \[
    650-572=78.
    \]
    La réduction est de $78$ euros.
\end{enumerate}
\end{correctionbox}

\begin{correctionbox}{Agrandissement}
Le coefficient d'agrandissement est $2{,}4$.

\begin{enumerate}
    \item Une longueur de $5$ cm devient :
    \[
    5\times 2{,}4=12\text{ cm}.
    \]
    
    \item Une longueur de $12$ cm devient :
    \[
    12\times 2{,}4=28{,}8\text{ cm}.
    \]
    
    \item Si la longueur agrandie est $36$ cm, alors la longueur initiale vaut :
    \[
    \frac{36}{2{,}4}=15\text{ cm}.
    \]
\end{enumerate}
\end{correctionbox}

\subsection{Corrections -- Grandeurs}

\begin{correctionbox}{Volume d'une boule}
Le rayon est :
\[
r=6\text{ cm}.
\]

Le volume d'une boule est :
\[
V=\frac{4}{3}\pi r^3.
\]

Donc :
\[
V=\frac{4}{3}\pi\times 6^3.
\]

Or :
\[
6^3=216.
\]

Donc :
\[
V=\frac{4}{3}\pi\times 216.
\]

Comme :
\[
\frac{216}{3}=72,
\]
on obtient :
\[
V=4\times 72\pi=288\pi.
\]

Le volume exact est :
\[
288\pi\text{ cm}^3.
\]

Une valeur approchée est :
\[
288\pi\approx 904{,}8\text{ cm}^3.
\]
\end{correctionbox}

\begin{correctionbox}{Vitesse et distance}
La voiture roule à $90$ km/h.

\begin{enumerate}
    \item Pour convertir en m/s, on divise par $3{,}6$ :
    \[
    90\div 3{,}6=25.
    \]
    Donc :
    \[
    90\text{ km/h}=25\text{ m/s}.
    \]
    
    \item En $2$ secondes, elle parcourt :
    \[
    25\times 2=50\text{ m}.
    \]
    
    \item En $5$ secondes, elle parcourt :
    \[
    25\times 5=125\text{ m}.
    \]
\end{enumerate}
\end{correctionbox}

\begin{correctionbox}{Débit}
Le débit est $15$ L/min.

\begin{enumerate}
    \item En $8$ minutes, le volume écoulé est :
    \[
    V=15\times 8=120\text{ L}.
    \]
    
    \item Comme :
    \[
    1\text{ m}^3=1000\text{ L},
    \]
    on a :
    \[
    120\text{ L}=0{,}120\text{ m}^3.
    \]
    
    \item Pour remplir $450$ L :
    \[
    t=\frac{450}{15}=30.
    \]
    Il faut $30$ minutes.
\end{enumerate}
\end{correctionbox}

\newpage

%================================================
\section{Grand devoir bilan}

\begin{center}
{\Large \textbf{Devoir bilan -- Données, probabilités, proportionnalité et grandeurs}}\\
Durée conseillée : 1 h 30
\end{center}

\subsection*{Exercice 1 -- Statistiques et histogramme}

Une enquête est réalisée auprès de $200$ élèves sur leur temps quotidien de trajet entre leur domicile et le collège.

\[
\begin{array}{|c|c|}
\hline
\text{Temps en minutes} & \text{Effectif} \\
\hline
\intervalle{0}{10} & 24 \\
\intervalle{10}{20} & 56 \\
\intervalle{20}{30} & 72 \\
\intervalle{30}{40} & 32 \\
\intervalle{40}{50} & 16 \\
\hline
\end{array}
\]

\begin{enumerate}
    \item Vérifier l'effectif total.
    \item Calculer la fréquence des élèves ayant un trajet entre $20$ et $30$ minutes.
    \item Calculer la fréquence des élèves ayant un trajet strictement inférieur à $20$ minutes.
    \item Représenter cette série par un histogramme.
    \item Donner une interprétation de la classe la plus représentée.
    \item Donner une estimation de l'étendue à partir des bornes des classes.
\end{enumerate}

\subsection*{Exercice 2 -- Probabilités}

Une urne contient $2$ boules rouges, $3$ boules bleues et $1$ boule verte.

On tire une boule au hasard, on la remet dans l'urne, puis on tire une deuxième boule.

\begin{enumerate}
    \item Combien y a-t-il de boules dans l'urne ?
    \item Quelle est la probabilité de tirer une boule rouge au premier tirage ?
    \item Construire un raisonnement permettant de calculer la probabilité de tirer deux boules bleues.
    \item Calculer la probabilité de ne pas tirer de boule verte au premier tirage.
    \item Expliquer pourquoi, si l'on répète l'expérience un très grand nombre de fois, la fréquence d'apparition d'une boule verte au premier tirage devrait se stabiliser autour de $\frac{1}{6}$.
\end{enumerate}

\subsection*{Exercice 3 -- Proportionnalité et pourcentages}

Un magasin vend une tablette au prix initial de $240$ euros.

Pendant les soldes, son prix baisse de $15\%$.

Une semaine plus tard, le nouveau prix augmente de $10\%$.

\begin{enumerate}
    \item Calculer le prix après la baisse de $15\%$.
    \item Calculer le prix après l'augmentation de $10\%$.
    \item Le prix final est-il égal au prix initial ? Justifier.
    \item Donner le coefficient multiplicateur global.
    \item En déduire le pourcentage global d'évolution.
\end{enumerate}

\subsection*{Exercice 4 -- Grandeurs et volumes}

Un réservoir est formé d'un cylindre de rayon $5$ cm et de hauteur $20$ cm, surmonté d'une demi-boule de même rayon.

\begin{enumerate}
    \item Calculer le volume du cylindre.
    \item Calculer le volume de la demi-boule.
    \item Calculer le volume total exact.
    \item Donner une valeur approchée au cm$^3$ près.
    \item Convertir ce volume en litres.
\end{enumerate}

\newpage

%================================================
\section{Correction détaillée du devoir bilan}

\subsection*{Correction de l'exercice 1}

\begin{enumerate}
    \item L'effectif total est :
    \[
    24+56+72+32+16=200.
    \]
    L'effectif total est bien $200$.
    
    \item La fréquence des élèves ayant un trajet entre $20$ et $30$ minutes est :
    \[
    f=\frac{72}{200}=0{,}36=36\%.
    \]
    
    \item Les élèves ayant un trajet strictement inférieur à $20$ minutes sont dans les classes :
    \[
    \intervalle{0}{10}\quad \text{et}\quad \intervalle{10}{20}.
    \]
    Leur effectif est :
    \[
    24+56=80.
    \]
    La fréquence est :
    \[
    \frac{80}{200}=0{,}4=40\%.
    \]
    
    \item Les classes ont toutes une amplitude de $10$ minutes. On peut donc tracer un histogramme avec les effectifs comme hauteurs.
    
    \item La classe la plus représentée est $\intervalle{20}{30}$ avec $72$ élèves. Cela signifie que le temps de trajet le plus fréquent dans cette enquête est compris entre $20$ et $30$ minutes.
    
    \item Les données vont de $0$ à moins de $50$ minutes. Une estimation de l'étendue est donc :
    \[
    50-0=50\text{ minutes}.
    \]
\end{enumerate}

\subsection*{Correction de l'exercice 2}

\begin{enumerate}
    \item Il y a :
    \[
    2+3+1=6
    \]
    boules dans l'urne.
    
    \item La probabilité de tirer une boule rouge au premier tirage est :
    \[
    \frac{2}{6}=\frac{1}{3}.
    \]
    
    \item Comme le tirage est fait avec remise, la composition de l'urne ne change pas.
    
    La probabilité de tirer une boule bleue à chaque tirage est :
    \[
    \frac{3}{6}=\frac{1}{2}.
    \]
    
    Pour tirer deux boules bleues :
    \[
    \frac{1}{2}\times \frac{1}{2}=\frac{1}{4}.
    \]
    
    On peut aussi raisonner par dénombrement : il y a $6\times 6=36$ couples possibles de boules distinguées. Les boules bleues sont au nombre de $3$, donc il y a $3\times 3=9$ couples favorables.
    \[
    \Prob(\text{deux bleues})=\frac{9}{36}=\frac{1}{4}.
    \]
    
    \item La probabilité de tirer une boule verte est :
    \[
    \frac{1}{6}.
    \]
    Donc la probabilité de ne pas tirer une boule verte est :
    \[
    1-\frac{1}{6}=\frac{5}{6}.
    \]
    
    \item La probabilité théorique de tirer une boule verte au premier tirage est :
    \[
    \frac{1}{6}.
    \]
    Si on répète l'expérience un très grand nombre de fois, la fréquence observée d'apparition d'une boule verte devrait se stabiliser autour de cette probabilité.
\end{enumerate}

\subsection*{Correction de l'exercice 3}

\begin{enumerate}
    \item Une baisse de $15\%$ correspond au coefficient :
    \[
    1-\frac{15}{100}=0{,}85.
    \]
    Le prix après baisse est :
    \[
    240\times 0{,}85=204.
    \]
    Le prix devient $204$ euros.
    
    \item Une augmentation de $10\%$ correspond au coefficient :
    \[
    1+\frac{10}{100}=1{,}10.
    \]
    Le prix final est :
    \[
    204\times 1{,}10=224{,}40.
    \]
    
    \item Le prix final est $224{,}40$ euros, alors que le prix initial était $240$ euros. Ils ne sont donc pas égaux.
    
    \item Le coefficient multiplicateur global est :
    \[
    0{,}85\times 1{,}10=0{,}935.
    \]
    
    \item Le coefficient $0{,}935$ correspond à :
    \[
    1-0{,}935=0{,}065.
    \]
    Donc l'évolution globale est une baisse de :
    \[
    6{,}5\%.
    \]
\end{enumerate}

\subsection*{Correction de l'exercice 4}

Le rayon est :
\[
r=5\text{ cm}.
\]

La hauteur du cylindre est :
\[
h=20\text{ cm}.
\]

\begin{enumerate}
    \item Volume du cylindre :
    \[
    V_1=\pi r^2h=\pi\times 5^2\times 20.
    \]
    Donc :
    \[
    V_1=500\pi\text{ cm}^3.
    \]
    
    \item Volume d'une boule complète :
    \[
    V_{\text{boule}}=\frac{4}{3}\pi r^3.
    \]
    Donc :
    \[
    V_{\text{boule}}=\frac{4}{3}\pi\times 5^3
    =\frac{4}{3}\pi\times 125
    =\frac{500}{3}\pi.
    \]
    Le volume d'une demi-boule est :
    \[
    V_2=\frac{1}{2}\times \frac{500}{3}\pi=\frac{250}{3}\pi.
    \]
    
    \item Le volume total exact est :
    \[
    V=500\pi+\frac{250}{3}\pi.
    \]
    On met au même dénominateur :
    \[
    500\pi=\frac{1500}{3}\pi.
    \]
    Donc :
    \[
    V=\frac{1500}{3}\pi+\frac{250}{3}\pi=\frac{1750}{3}\pi.
    \]
    
    \item Valeur approchée :
    \[
    V=\frac{1750}{3}\pi\approx 1832{,}6.
    \]
    Au cm$^3$ près :
    \[
    V\approx 1833\text{ cm}^3.
    \]
    
    \item Comme :
    \[
    1\text{ L}=1000\text{ cm}^3,
    \]
    on obtient :
    \[
    1833\text{ cm}^3\approx 1{,}833\text{ L}.
    \]
\end{enumerate}

\newpage

%================================================
\section{Fiche de synthèse finale}

\begin{retenirbox}{Statistiques}
\begin{itemize}
    \item L'effectif total est le nombre total de données.
    \item L'effectif d'une valeur est le nombre de fois où cette valeur apparaît.
    \item La fréquence est :
    \[
    f=\frac{\text{effectif}}{\text{effectif total}}.
    \]
    \item La somme des fréquences vaut $1$, c'est-à-dire $100\%$.
    \item L'étendue est :
    \[
    \text{max}-\text{min}.
    \]
    \item Un diagramme en bâtons représente des valeurs séparées.
    \item Un histogramme représente des classes d'intervalles.
\end{itemize}
\end{retenirbox}

\begin{retenirbox}{Probabilités}
\begin{itemize}
    \item Une probabilité est comprise entre $0$ et $1$.
    \item En équiprobabilité :
    \[
    \Prob(A)=\frac{\text{nombre d'issues favorables}}{\text{nombre total d'issues}}.
    \]
    \item Pour l'événement contraire :
    \[
    \Prob(\overline{A})=1-\Prob(A).
    \]
    \item Pour une expérience à deux épreuves, on peut utiliser un tableau à double entrée ou un arbre.
    \item Lorsque le nombre d'essais augmente, les fréquences observées se stabilisent autour des probabilités théoriques.
\end{itemize}
\end{retenirbox}

\begin{retenirbox}{Proportionnalité}
\begin{itemize}
    \item Deux grandeurs sont proportionnelles si l'une s'obtient en multipliant l'autre par un coefficient constant.
    \item Une situation de proportionnalité se modélise par :
    \[
    f(x)=ax.
    \]
    \item Augmenter de $p\%$ revient à multiplier par :
    \[
    1+\frac{p}{100}.
    \]
    \item Diminuer de $p\%$ revient à multiplier par :
    \[
    1-\frac{p}{100}.
    \]
\end{itemize}
\end{retenirbox}

\begin{retenirbox}{Grandeurs}
\begin{itemize}
    \item Volume d'une boule :
    \[
    V=\frac{4}{3}\pi r^3.
    \]
    \item Volume d'un cylindre :
    \[
    V=\pi r^2h.
    \]
    \item Vitesse :
    \[
    v=\frac{d}{t}.
    \]
    \item Débit :
    \[
    D=\frac{V}{t}.
    \]
    \item Conversions importantes :
    \[
    1\text{ h}=3600\text{ s},
    \qquad
    1\text{ m}^3=1000\text{ L},
    \qquad
    1\text{ km/h}=\frac{1}{3{,}6}\text{ m/s}.
    \]
\end{itemize}
\end{retenirbox}

\begin{warningbox}{Derniers points de vigilance}
\begin{itemize}
    \item Toujours distinguer effectif et fréquence.
    \item Ne pas confondre histogramme et diagramme en bâtons.
    \item Toujours vérifier que les fréquences totalisent $1$ ou $100\%$.
    \item Toujours lister clairement les issues en probabilités.
    \item Toujours convertir les unités avant de calculer.
    \item Toujours donner l'unité finale d'une grandeur.
    \item Une baisse de $20\%$ suivie d'une hausse de $20\%$ ne ramène pas à la valeur initiale.
\end{itemize}
\end{warningbox}

\end{document}